\chapter{Triangulacje i grafy dualne kompleksów symplicjalnych}
\label{ch:tg-triangulacje}

Każdy skończony graf prosty można traktować jako kompleks
symplicjalny wymiaru co najwyżej jeden. Graf bez krawędzi
o niepustym zbiorze wierzchołków ma wymiar zero, a graf
zerowy ma wymiar $-1$. Grafy z pętlami lub krawędziami
równoległymi wymagają opisu komórkowego albo odpowiedniego
podziału krawędzi. Aby opisywać
powierzchnie i rozmaitości, potrzebujemy także dwuwymiarowych
trójkątów i sympleksów wyższych wymiarów. Zaczniemy od definicji,
a następnie sprawdzimy, jakie dodatkowe dane musi zawierać
triangulacja, żeby wyznaczać powierzchnię, a nie tylko jej graf.

\section{Kompleks symplicjalny i realizacja geometryczna}
\label{sec:tg-kompleks-symplicjalny}

\begin{definicja}[Abstrakcyjny kompleks symplicjalny]
Niech $V$ będzie skończonym zbiorem. \emph{Abstrakcyjnym
kompleksem symplicjalnym} na $V$ nazywamy rodzinę $K$
niepustych, skończonych podzbiorów $V$ taką, że jeżeli
$\sigma\in K$ i $\varnothing\ne\tau\subseteq\sigma$,
to $\tau\in K$. Zakładamy, że $\{v\}\in K$ dla każdego $v\in V$. Element $\sigma\in K$
o $p+1$ wierzchołkach nazywamy \emph{$p$-sympleksem}.
\emph{Wymiarem} $K$ jest
$\dim K:=\max\{|\sigma|-1:\sigma\in K\}$ dla $K\ne\varnothing$,
a dla pustego kompleksu przyjmujemy $\dim K=-1$.
\end{definicja}

Wierzchołki są zatem sympleksami wymiaru zero, krawędzie
sympleksami wymiaru jeden, a trójkąty sympleksami wymiaru dwa.
Warunek domknięcia na niepuste podzbiory mówi, że jeśli
trójkąt należy do kompleksu, to należą do niego także jego
trzy krawędzie i trzy wierzchołki. Nie utożsamiamy
trójkąta z samym jego brzegiem.

\begin{definicja}[Realizacja geometryczna]
Niech $K$ będzie skończonym kompleksem na zbiorze $V$.
Jego \emph{realizacją geometryczną} nazywamy
\[
|K|=\left\{(t_v)_{v\in V}\in\mathbb R^V:\
t_v\ge0,\quad\sum_{v\in V}t_v=1,\quad
\{v:t_v>0\}\in K\right\}.
\]
Nadajemy mu topologię podprzestrzeni $\mathbb R^V$.
Jest zwarty i Hausdorffa: stanowi skończoną sumę
domkniętych sympleksów, z których każdy jest zwarty.
Dla $K=\varnothing$ realizacja jest pusta.
Podzbiór, w którym $t_v=0$ dla $v\notin\sigma$, jest
geometrycznym sympleksem odpowiadającym $\sigma\in K$.
\end{definicja}

\begin{przyklad}[Trójkąt i jego brzeg]
Niech $V=\{a,b,c\}$. Kompleks zawierający wszystkie niepuste
podzbiory $V$ realizuje pełny trójkąt, homeomorficzny z dyskiem.
Jeżeli usuniemy jedyny dwusympleks $\{a,b,c\}$, pozostawiając
wierzchołki i trzy krawędzie, otrzymamy brzeg trójkąta,
homeomorficzny z $S^1$. Obie realizacje mają identyczny
jednowymiarowy szkielet, lecz różne grupy fundamentalne
i homologie. Sam graf nie odróżnia zatem tych obiektów.
\end{przyklad}

\begin{samepage}
\begin{center}
\begin{minipage}{.96\linewidth}\centering
\begin{tikzpicture}[scale=1.05,line cap=round]
\begin{scope}[shift={(-2.8,0)}]
 \coordinate (a) at (-1,-.7);\coordinate (b) at (1,-.7);
 \coordinate (c) at (0,1.05);
 \fill[tgBlue!15] (a)--(b)--(c)--cycle;
 \draw[tgEdge] (a)--(b)--(c)--cycle;
 \foreach \p in {a,b,c} \node[tgBlueVertex] at (\p) {};
 \node[text=tgNavy,font=\small] at (0,-1.25) {$|K|\cong D^2$};
\end{scope}
\begin{scope}[shift={(2.8,0)}]
 \coordinate (a) at (-1,-.7);\coordinate (b) at (1,-.7);
 \coordinate (c) at (0,1.05);
 \draw[tgAccent] (a)--(b)--(c)--cycle;
 \foreach \p in {a,b,c} \node[tgRedVertex] at (\p) {};
 \node[text=tgNavy,font=\small] at (0,-1.25) {$|\partial K|\cong S^1$};
\end{scope}
\end{tikzpicture}\par\smallskip

{\footnotesize Ten sam graf krawędziowy: po lewej z dołączoną
2-komórką, po prawej bez niej.}
\end{minipage}
\end{center}
\end{samepage}

\begin{definicja}[Szkielet i ogniwo]
Dla $p\ge0$ przez $K^{(p)}$ oznaczamy podkompleks
złożony ze wszystkich sympleksów wymiaru co najwyżej $p$.
Dla sympleksu $\sigma\in K$ definiujemy jego
\emph{ogniwo} (ang. \emph{link}) przez
$\operatorname{Lk}_K(\sigma)
=\{\varnothing\ne\tau\subseteq V\setminus\sigma:
\ \tau\cup\sigma\in K\}$.
Zbiorem wierzchołków ogniwa jest suma jego sympleksów,
niekoniecznie całe $V\setminus\sigma$. \emph{Gwiazdą domkniętą}
$\operatorname{St}_K(\sigma)$ nazywamy podkompleks
wszystkich niepustych ścian sympleksów zawierających $\sigma$.
\end{definicja}

Dla wierzchołka $v$ o niepustym ogniwie jego gwiazda
domknięta jest stożkiem nad ogniwem. Istotnie, punkt
gwiazdy różny od $v$ można jednoznacznie zapisać jako
$t v+(1-t)y$, gdzie $0\le t<1$, $y\in|\operatorname{Lk}_K(v)|$;
we współrzędnych barycentrycznych $t=t_v$, a pozostałe
współrzędne dzielimy przez $1-t$. Przy $t=1$ wszystkie
punkty podstawy zostają utożsamione z $v$. Otrzymana
ciągła bijekcja ze zwartego stożka do przestrzeni Hausdorffa
jest homeomorfizmem. Dla izolowanego $v$ gwiazda jest
jednym punktem. Zbiór $t_v>0$ jest otwarty w $|K|$,
a zbiory $t_v>1-\varepsilon$ dla $0<\varepsilon<1$
tworzą bazę otoczeń $v$ (wynika to z ograniczoności
pozostałych współrzędnych przez $\varepsilon$).
Po usunięciu $v$ każde z tych otoczeń ma postać
$|\operatorname{Lk}_K(v)|\times(1-\varepsilon,1)$.
Ten opis pozwoli rozpoznawać lokalną topologię powierzchni.

\section{Triangulowana powierzchnia i warunek ogniwa}
\label{sec:tg-link-manifold}

\begin{definicja}[Triangulacja]
Niech $M$ będzie przestrzenią topologiczną. \emph{Triangulacją}
$M$ nazywamy parę $(K,h)$, gdzie $K$ jest skończonym
kompleksem symplicjalnym, a $h:|K|\to M$ homeomorfizmem.
Ponieważ $K$ jest skończony, taka triangulacja może istnieć
wyłącznie dla zwartej $M$; w dalszych twierdzeniach rozważamy
powierzchnie zwarte i triangulacje skończone. Gdy $M$ jest
powierzchnią, wierzchołki, krawędzie i trójkąty $K$ nazywamy
elementami jej triangulacji. Triangulację
nazywamy \emph{czystą dwuwymiarową}, gdy $\dim K=2$
i każdy sympleks jest ścianą pewnego dwusympleksu
(dopuszczając sam dwusympleks jako jego ścianę).
\end{definicja}

Każda skończona triangulacja niepustej powierzchni bez
brzegu jest czysta dwuwymiarowa. Uzasadnienie korzysta
z lokalnej homologii singularnej: dla punktu $x$
we względnym wnętrzu maksymalnego $p$-sympleksu istnieje
otoczenie homeomorficzne z $\mathbb R^p$. Istotnie,
wewnętrzny punkt nie należy do innych sympleksów poza
tym i jego nadścianami, a maksymalność wyklucza nadściany;
skończoność pozwala odsunąć się od wszystkich pozostałych
domkniętych sympleksów. Przez wycięcie
$H_i(|K|,|K|\setminus\{x\};\mathbb Z)$ jest wówczas
równa $\mathbb Z$ dla $i=p$ i zeru w innych stopniach,
z ciągu dokładnego pary i retrakcji przebitego dysku
na sferę $S^{p-1}$ (dla $p=0$ para to punkt i zbiór pusty).
Na powierzchni bez brzegu ten sam rachunek w mapie
powierzchni daje jedyną niezerową grupę w stopniu $2$.
Zatem każdy maksymalny sympleks ma wymiar dwa.
Korzystamy tu z klasycznych twierdzeń homologii singularnej:
wycięcia, ciągu dokładnego pary i homologii sfer.

Nie każdy skończony kompleks złożony z trójkątów jest
powierzchnią. Dwa trójkąty sklejone tylko w jednym
wierzchołku mają w tym punkcie otoczenie niehomeomorficzne
z dyskiem. Potrzebujemy zatem lokalnego kryterium.

\begin{twierdzenie}[Kryterium ogniwa dla powierzchni zamkniętej]
\label{thm:tg-link-surface}
Niech $K$ będzie skończonym, spójnym, czystym
kompleksem symplicjalnym wymiaru dwa. Następujące warunki
są równoważne:
\begin{enumerate}[label=(\roman*)]
\item $|K|$ jest zamkniętą powierzchnią topologiczną;
\item dla każdego wierzchołka $v\in K$ realizacja
$|\operatorname{Lk}_K(v)|$ jest okręgiem.
\end{enumerate}
W szczególności, jeżeli zachodzą te warunki, każda
krawędź należy do dokładnie dwóch trójkątów.
\end{twierdzenie}

\begin{proof}
Załóżmy najpierw (ii). Gwiazda każdego wierzchołka jest
stożkiem nad okręgiem, a więc jest dyskiem; jego część
$t_v>0$ jest otwartym dyskiem wokół wierzchołka. Wnętrza trójkątów
mają otoczenia dyskowe wprost z geometrii sympleksu.
Każda krawędź należy do dwóch trójkątów: w ogniwie jednego
jej końca odpowiada jej wierzchołek, który w cyklu ma
dokładnie dwóch sąsiadów. Zatem również punkt we wnętrzu
krawędzi ma otoczenie powstałe z dwóch półdysków i jest
punktem wewnętrznym powierzchni. Wszystkie punkty mają
więc otoczenia homeomorficzne z $\mathbb R^2$. Ponieważ
$|K|$ jest zwarty, spójny, Hausdorffa i ma przeliczalną
bazę topologii jako podprzestrzeń skończenie wymiarowej
przestrzeni euklidesowej, jest zamkniętą powierzchnią.

W drugą stronę załóżmy (i). Dla punktu $x$ we wnętrzu
krawędzi $e$ małe otoczenie w $|K|$ jest sklejeniem
$k$ półdysków wzdłuż ich średnic, gdzie $k$ jest liczbą
trójkątów zawierających $e$. Ogniwo krawędzi ma $k$
punktów. Małe otwarte otoczenie $U$ punktu $x$ jest ściągalne:
w każdym półdysku wykonujemy radialne ściąganie do $x$,
zgodne na wspólnej średnicy. Po usunięciu $x$ radialna
deformacja na ustaloną mniejszą odległość od $x$ daje
graf z dwoma wierzchołkami i $k$ łukami pomiędzy nimi,
po jednym z każdego półdysku. Ponieważ $k\ge1$ z czystości
kompleksu, graf ten jest spójny i jego pierwsza homologia
jest $\mathbb Z^{k-1}$: wybieramy jeden łuk jako drzewo,
a każdy z pozostałych zamyka niezależny cykl.
Z wycięcia i ciągu dokładnego pary otrzymujemy
\[
H_2(|K|,|K|\setminus\{x\};\mathbb Z)
\cong H_2(U,U\setminus\{x\};\mathbb Z)
\cong H_1(U\setminus\{x\};\mathbb Z)
\cong \mathbb Z^{k-1}.
\]
Drugi izomorfizm wynika z $H_2(U)=H_1(U)=0$.
Dla punktu powierzchni bez brzegu ta grupa jest
izomorficzna z $\mathbb Z$, zatem $k=2$.

Ogniwo wierzchołka $v$ jest skończonym grafem.
Jego wierzchołkowi $u$ odpowiada krawędź $\{v,u\}$,
a jego krawędziom incydentnym z $u$ trójkąty zawierające
$\{v,u\}$. Każdy wierzchołek ogniwa ma więc stopień dwa,
czyli ogniwo jest sumą rozłączną cykli. Niech $U$ będzie
małym otoczeniem stożkowym $v$ opisanym powyżej.
Jest ściągalne, a $U\setminus\{v\}$ deformacyjnie
retrahuje się na ogniwo. Wycięcie i ciąg dokładny
pary w stopniu jeden dają
\[
 H_1(|K|,|K|\setminus\{v\};\mathbb Z)
 \cong \widetilde H_0(|\operatorname{Lk}_K(v)|;\mathbb Z).
\]
Lewa strona jest zerowa dla punktu powierzchni bez brzegu.
Prawa ma rangę o jeden mniejszą niż liczba składowych
niepustego ogniwa. Zatem ogniwo jest spójne, a więc
jest pojedynczym cyklem. Ta argumentacja nie zakłada,
że dowolne małe otoczenie wierzchołka jest samo dyskiem.
\end{proof}

\begin{przyklad}[Brzeg czworościanu]
Niech $K$ składa się z czterech trójkątów, sześciu krawędzi
i czterech wierzchołków tworzących brzeg czworościanu.
Każde ogniwo wierzchołka jest trójkątnym cyklem, więc
$|K|$ jest zamkniętą powierzchnią. Mamy
$\chi(|K|)=4-6+4=2$, a ponieważ $|K|$ jest homeomorficzny z brzegiem
wypukłego czworościanu w $\mathbb R^3$, rzut radialny
z punktu jego wnętrza na sferę jednostkową daje
homeomorfizm z $S^2$: każdy promień przecina brzeg
w dokładnie jednym punkcie, a ciągła bijekcja
ze zwartego brzegu do sfery jest homeomorfizmem.
W tym przykładzie cztery dwusympleksy odpowiadają również
czterem ścianom standardowego osadzenia szkieletu
czworościanu na sferze. W ogólności trzeba jednak odróżniać
dwusympleksy konkretnego kompleksu od ścian dowolnego
osadzenia samego grafu.
\end{przyklad}

\begin{samepage}
\begin{center}
\begin{minipage}{.96\linewidth}\centering
\begin{tikzpicture}[scale=1.05,line join=round,line cap=round]
  \coordinate (a) at (-1.3,-.8);\coordinate (b) at (1.3,-.8);
  \coordinate (c) at (.35,1.4);\coordinate (d) at (-.22,.13);
  \fill[tgBlue!11,rounded corners=5pt] (-2,-1.02) rectangle (2,1.92);
  \fill[tgCoral!18] (a)--(b)--(d)--cycle;
  \fill[tgTeal!18] (a)--(c)--(d)--cycle;
  \fill[tgGold!20] (b)--(c)--(d)--cycle;
  \draw[tgEdge] (a)--(b)--(c)--cycle;
  \draw[tgEdge] (a)--(d)--(b) (c)--(d);
  \node[tgBlueVertex,label=below left:$a$] at (a) {};
  \node[tgBlueVertex,label=below right:$b$] at (b) {};
  \node[tgBlueVertex,label=above right:$c$] at (c) {};
  \node[tgBlueVertex,label=below:$d$] at (d) {};
  \node[font=\scriptsize,text=tgNavy] at (-.65,1.67) {ściana zewnętrzna $abc$};
  \node[font=\small,text=tgNavy] at (0,-1.38) {$V=4,\ E=6,\ F=4,\ \chi=2$};
\end{tikzpicture}\par\smallskip

{\footnotesize Schemat kombinatoryczny powierzchni czworościanu;
trzy ściany wewnętrzne i ściana zewnętrzna (niebieskie tło)
odpowiadają czterem trójkątom. Obramowanie tła nie jest krawędzią grafu.}
\end{minipage}
\end{center}
\end{samepage}

\section{Łańcuchy symplicjalne i związek z homologią}
\label{sec:tg-simplicial-chain}

\begin{definicja}[Operator brzegu]
Dla skończonego kompleksu symplicjalnego $K$ niech
$C_p(K;\mathbb Z)$ będzie wolną grupą abelową generowaną
przez zorientowane $p$-sympleksy, z relacją
$[v_{\pi(0)},\ldots,v_{\pi(p)}]
=\operatorname{sgn}(\pi)[v_0,\ldots,v_p]$.
Dla $p\ge1$ definiujemy
\[
\partial_p[v_0,\ldots,v_p]
=\sum_{i=0}^p(-1)^i
[v_0,\ldots,\widehat v_i,\ldots,v_p],
\]
a $\partial_0=0$. Wzór jest zgodny z relacją orientacji:
wystarczy sprawdzić zamianę dwóch sąsiednich wierzchołków
o indeksach $i,i+1$, ponieważ takie transpozycje generują
wszystkie permutacje. Składniki, w których usuwamy inny
wierzchołek, zmieniają znak przez transpozycję w zachowanej
ścianie. Dwa składniki, w których usuwamy jeden z
zamienianych wierzchołków, zamieniają się miejscami
i otrzymują przeciwne znaki $(-1)^i$ oraz $(-1)^{i+1}$.
Cały brzeg zmienia więc znak, zgodnie z relacją.
Homologie symplicjalne są
$H_p(K;\mathbb Z):=\ker\partial_p/\operatorname{im}\partial_{p+1}$.
Dla przemiennego pierścienia $R$ z jedynką przyjmujemy
$C_p(K;R)=C_p(K;\mathbb Z)\otimes_{\mathbb Z}R$
i rozszerzamy operatory brzegowe $R$-liniowo.
Homologie z tymi współczynnikami definiujemy jako
jądro modulo obraz nowych operatorów. Przyjmujemy
$C_{-1}=0$; używamy homologii niezredukowanej.
Korzystamy z klasycznego twierdzenia o zgodności homologii
symplicjalnej i singularnej realizacji $|K|$.
\end{definicja}

\begin{stwierdzenie}[Brzeg brzegu]
Dla każdego skończonego kompleksu symplicjalnego $K$
i każdego $p\ge1$ zachodzi $\partial_{p-1}\partial_p=0$.
\end{stwierdzenie}

\begin{proof}
Dla $p=1$ teza wynika z $\partial_0=0$. Dla $p\ge2$
w rozwinięciu $\partial_{p-1}\partial_p$ każda ściana
powstała przez usunięcie dwóch różnych wierzchołków
pojawia się dwukrotnie. Jeśli najpierw usuniemy wierzchołek
o indeksie $i$, a następnie o indeksie $j>i$, znak jest
$(-1)^i(-1)^{j-1}$. Przy odwrotnej kolejności znak jest
$(-1)^j(-1)^i$. Ich suma wynosi zero. Sumując po wszystkich
parach $i<j$, otrzymujemy równość operatorową.
\end{proof}

\begin{przyklad}[Co robi dołączenie trójkąta]
Dla brzegu trójkąta o wierzchołkach $0,1,2$ weźmy łańcuch
$a[0,1]+b[1,2]+c[0,2]$. Jego brzeg to
$(-a-c)[0]+(a-b)[1]+(b+c)[2]$.
Warunek zerowego brzegu daje $a=b=-c$, więc jądro
$\partial_1$ jest generowane swobodnie przez
$z=[0,1]+[1,2]-[0,2]$. Ponieważ $C_2=0$, otrzymujemy
$H_1(\partial\Delta^2;\mathbb Z)\cong\mathbb Z$.
Po dołączeniu trójkąta $[0,1,2]$ otrzymujemy
$z=\partial_2[0,1,2]$, więc klasa $z$ zanika i
$H_1(\Delta^2;\mathbb Z)=0$. W obu przypadkach
obraz $\partial_1$ jest generowany przez $[1]-[0]$
i $[2]-[0]$: trzeci brzeg $[2]-[1]$ jest ich różnicą.
W ilorazie $C_0/\operatorname{im}\partial_1$ wszystkie
trzy generatory reprezentują więc tę samą klasę, a suma
współczynników łańcucha wyznacza izomorfizm z $\mathbb Z$.
Zatem $H_0\cong\mathbb Z$ w obu przypadkach. To bezpośredni rachunek ilustrujący
różnicę między grafem i wypełnionym kompleksem.
\end{przyklad}

\begin{stwierdzenie}[Charakterystyka Eulera z homologii]
Niech $K$ będzie skończonym kompleksem symplicjalnym,
a $f_p$ liczbą jego $p$-sympleksów. Niech $\mathbb F$
będzie dowolnym ciałem. Jeżeli
$b_p=\dim_{\mathbb F}H_p(K;\mathbb F)$, to
\[
\sum_{p\ge0}(-1)^pf_p=\sum_{p\ge0}(-1)^pb_p.
\]
\end{stwierdzenie}

\begin{proof}
Niech $r_p$ oznacza rząd operatora
$\partial_p:C_p(K;\mathbb F)\to C_{p-1}(K;\mathbb F)$;
przyjmujemy $r_0=0$. Z twierdzenia o rzędzie i jądrze
oraz definicji homologii mamy
$b_p=(f_p-r_p)-r_{p+1}$, a więc
$f_p=b_p+r_p+r_{p+1}$. Po wymnożeniu przez $(-1)^p$
i zsumowaniu wkłady $r_p$ znoszą się parami, pozostawiając
sumę liczb Bettiego. Wszystkie sumy są skończone.
\end{proof}

\section{Graf dualny triangulacji i dane sklejania}
\label{sec:tg-dual-triangulation}

\begin{definicja}[Graf dualny triangulacji powierzchni]
Niech $K$ będzie skończoną triangulacją spójnej,
zamkniętej powierzchni. \emph{Graf dualny triangulacji}
ma po jednym wierzchołku dla każdego trójkąta $K$
i po jednej krawędzi dla każdej krawędzi $K$, łącząc
wierzchołki dwóch trójkątów incydentnych z daną krawędzią.
\end{definicja}

\begin{stwierdzenie}
Dla triangulacji spójnej powierzchni zamkniętej graf dualny jest
prosty, spójny i trójregularny, czyli każdy jego wierzchołek ma
stopień trzy. Jeżeli triangulacja ma
$f_2$ trójkątów i $f_1$ krawędzi, to
$3f_2=2f_1$.
\end{stwierdzenie}

\begin{proof}
Każdy trójkąt ma trzy boki, a każda krawędź zamkniętej
powierzchni jest wspólna dokładnie dwóm trójkątom.
Stąd stopień dualnego wierzchołka jest równy trzy,
a podwójne zliczenie par (trójkąt, jego bok) daje
$3f_2=2f_1$. Nie ma pętli, ponieważ krawędź jest
wspólna dwóm różnym trójkątom. Dwa różne trójkąty
abstrakcyjnego kompleksu nie mogą mieć dwóch wspólnych
boków: ich suma wierzchołków wyznaczałaby ten sam
trójelementowy zbiór. Nie ma więc krawędzi równoległych.

Dla ustalonego wierzchołka triangulacji $v$ trójkąty
zawierające $v$ odpowiadają krawędziom jego ogniwa,
które na mocy twierdzenia \ref{thm:tg-link-surface}
jest cyklem. Kolejne krawędzie cyklu mają wspólny
wierzchołek $u$, czyli odpowiadające trójkąty mają
wspólny bok $\{v,u\}$. Wszystkie trójkąty zawierające
$v$ należą zatem do jednej składowej grafu dualnego.
Dla każdej składowej grafu dualnego bierzemy sumę
domkniętych trójkątów do niej należących. Sumy te są
rozłączne: jeżeli dwa trójkąty się przecinają, mają
wspólny wierzchołek i należą do tej samej składowej.
Są też domknięte i pokrywają $|K|$ z czystości kompleksu.
Gdyby składowych było więcej niż jedna, otrzymalibyśmy
rozbicie spójnej powierzchni na rozłączne niepuste
zbiory domknięte i otwarte (dopełnienie każdej sumy
jest skończoną sumą domkniętych trójkątów).
Stąd graf dualny jest spójny.
\end{proof}

\begin{uwaga}[Dlaczego sam graf dualny nie wystarczy]
Graf dualny mówi, które trójkąty sąsiadują, ale nie pamięta,
\emph{które ich boki} są sklejane, w jakiej kolejności i czy
sklejenie odwraca orientację. Rozkłady powierzchni za pomocą
trójkątów z identyfikowanymi bokami bywają wygodniej opisane
przez \emph{kompleksy typu $\Delta$}, które nie muszą być
abstrakcyjnymi kompleksami symplicjalnymi. Dla przykładu
sferę można uzyskać przez sklejenie dwóch trójkątów wzdłuż
odpowiadających boków, a torus otrzymujemy z kwadratu z jedną przekątną,
identyfikując przeciwległe boki kwadratu przez przesunięcia.
W pierwszym modelu mamy trzy wierzchołki i trzy krawędzie,
w drugim wszystkie cztery narożniki dają jeden wierzchołek,
dwie pary boków i przekątna dają trzy krawędzie. W obu
modelach graf dualny jest grafem dwóch wierzchołków
i trzech krawędzi równoległych, mimo że
$\chi(S^2)=2$, a $\chi(T^2)=0$. Rozróżniają je dopiero
dane sklejeń; są to modele typu $\Delta$, nie symplicjalne
triangulacje z dwiema dwukomórkami.
\end{uwaga}

\section{Ruchy Pachnera na triangulacjach powierzchni}
\label{sec:tg-pachner}

\begin{definicja}[Lokalne ruchy bistellarne w wymiarze dwa]
Niech $K$ będzie triangulacją zamkniętej powierzchni.
Ruch $1\leftrightarrow3$ zastępuje pojedynczy trójkąt
trzema trójkątami przez wprowadzenie nowego wierzchołka
w jego wnętrzu. Ruch odwrotny $3\to1$ jest dopuszczalny,
gdy ogniwo usuwanego wierzchołka $x$ jest cyklem o trzech
wierzchołkach $a,b,c$ i trójkąt $\{a,b,c\}$ nie należy
do kompleksu. Usuwamy wówczas gwiazdę $x$ z wyjątkiem
jej brzegu i wstawiamy trójkąt $\{a,b,c\}$.
Warunek braku tego trójkąta jest konieczny, np. ruch
nie jest dopuszczalny na brzegu czworościanu.
Ruch $2\leftrightarrow2$ usuwa wspólną przekątną dwóch
różnych trójkątów, których suma jest podkompleksem
homeomorficznym z domkniętym dyskiem o brzegu złożonym
z czterech różnych krawędzi, po czym wprowadza drugą
przekątną tego czworokąta. Wymagamy czterech różnych wierzchołków brzegu i tego,
by nowa przekątna nie należała do kompleksu. Wtedy
nowe trójkąty także nie należą do niego, ponieważ
zawierałyby tę przekątną. Warunki te zapewniają,
że nowe sympleksy spotykają pozostałą część kompleksu
tylko w ścianach brzegu wymienianego dysku.
\end{definicja}

\begin{samepage}
\begin{center}
\begin{minipage}{.96\linewidth}\centering
\begin{tikzpicture}[scale=.94,line cap=round,line join=round,>=Stealth]
 \begin{scope}[shift={(-3,0)}]
  \coordinate (a) at (-1,-.7);\coordinate (b) at (1,-.7);
  \coordinate (c) at (0,1.0);\coordinate (o) at (0,-.1);
  \fill[tgBlue!12] (a)--(b)--(c)--cycle;
  \draw[tgEdge] (a)--(b)--(c)--cycle;
  \node[tgBlueVertex] at (a) {};\node[tgBlueVertex] at (b) {};\node[tgBlueVertex] at (c) {};
  \node[text=tgNavy,font=\small] at (0,-1.2) {$1$ trójkąt};
 \end{scope}
 \draw[tgTeal,-{Stealth[length=3mm]},line width=1.5pt] (-1.45,0)--(-.75,0);
 \begin{scope}[shift={(1,0)}]
  \coordinate (a) at (-1,-.7);\coordinate (b) at (1,-.7);
  \coordinate (c) at (0,1.0);\coordinate (o) at (0,-.1);
  \fill[tgTeal!12] (a)--(b)--(o)--cycle;
  \fill[tgBlue!12] (a)--(c)--(o)--cycle;
  \fill[tgGold!16] (b)--(c)--(o)--cycle;
  \draw[tgEdge] (a)--(b)--(c)--cycle;
  \draw[tgAccent] (a)--(o)--(b) (c)--(o);
  \foreach \p in {a,b,c,o} \node[tgBlueVertex] at (\p) {};
  \node[text=tgNavy,font=\small] at (0,-1.2) {$3$ trójkąty};
 \end{scope}
\end{tikzpicture}\par\smallskip

{\footnotesize Ruch $1\to3$: powierzchnia oraz brzeg modyfikowanego dysku pozostają te same;
zmienia się tylko podział symplicjalny dysku.}
\end{minipage}
\end{center}
\end{samepage}

\begin{samepage}
\begin{center}
\begin{minipage}{.96\linewidth}\centering
\begin{tikzpicture}[scale=.84,line cap=round,line join=round,>=Stealth]
 \begin{scope}[shift={(-2.5,0)}]
  \coordinate (a) at (-1.25,-.85);\coordinate (b) at (1.25,-.85);
  \coordinate (c) at (1.05,.85);\coordinate (d) at (-1.05,.85);
  \fill[tgTeal!14] (a)--(b)--(c)--cycle;
  \fill[tgBlue!14] (a)--(c)--(d)--cycle;
  \draw[tgEdge] (a)--(b)--(c)--(d)--cycle;
  \draw[tgAccent] (a)--(c);
  \foreach \p in {a,b,c,d} \node[tgBlueVertex] at (\p) {};
  \node[font=\small,text=tgNavy] at (0,-1.30) {przekątna $ac$};
 \end{scope}
 \draw[tgTeal,-{Stealth[length=3mm]},line width=1.5pt] (-.64,0)--(.64,0);
 \begin{scope}[shift={(2.5,0)}]
  \coordinate (a) at (-1.25,-.85);\coordinate (b) at (1.25,-.85);
  \coordinate (c) at (1.05,.85);\coordinate (d) at (-1.05,.85);
  \fill[tgGold!15] (a)--(b)--(d)--cycle;
  \fill[tgCoral!12] (b)--(c)--(d)--cycle;
  \draw[tgEdge] (a)--(b)--(c)--(d)--cycle;
  \draw[tgAccent] (b)--(d);
  \foreach \p in {a,b,c,d} \node[tgRedVertex] at (\p) {};
  \node[font=\small,text=tgNavy] at (0,-1.30) {przekątna $bd$};
 \end{scope}
\end{tikzpicture}\par\smallskip

{\footnotesize Ruch $2\leftrightarrow2$: dwie triangulacje tego samego
czworokąta, z niezmienionym brzegiem.}
\end{minipage}
\end{center}
\end{samepage}

\begin{stwierdzenie}[Niezmienniczość topologii i charakterystyki Eulera]
Każdy dopuszczalny ruch $1\leftrightarrow3$ i
$2\leftrightarrow2$ zachowuje typ homeomorficzny
triangulowanej powierzchni oraz jej charakterystykę Eulera.
\end{stwierdzenie}

\begin{proof}
Obie operacje odbywają się wewnątrz domkniętego dysku
triangulacji i nie zmieniają sposobu sklejenia jego brzegu
z resztą powierzchni. Ruchy zastępują dysk inną
triangulacją tego samego dysku. Każdą z lokalnych
triangulacji można zrealizować w tym samym wypukłym
trójkącie lub czworokącie, z identycznymi położeniami
wierzchołków brzegu. Afiniczne odwzorowania poszczególnych
sympleksów w ten model są zgodne na wspólnych ścianach
i dają homeomorfizmy dysków, których złożenie zachowuje
brzeg punkt po punkcie. Przedłużając przez tożsamość
poza dyskiem, otrzymujemy homeomorfizm całej powierzchni.
W ruchu $1\to3$ przybywa jeden wierzchołek, trzy krawędzie
i dwie ściany dwuwymiarowe netto. Zmiana $f_0-f_1+f_2$
wynosi $1-3+2=0$. Ruch $2\to2$ nie zmienia żadnej z tych
liczb; odwrotne ruchy zachowują je również.
\end{proof}

\begin{stwierdzenie}[Subdywizja krawędzi przez dwa ruchy]
Niech $e=ab$ będzie krawędzią triangulacji zamkniętej
powierzchni, przylegającą do dwóch trójkątów $abc$
i $abd$ o różnych wierzchołkach $c,d$. Wtedy podział krawędzi
$ab$ przez nowy wierzchołek
jest kombinatorycznie równoważny wykonaniu jednego
ruchu $1\to3$ oraz jednego ruchu $2\to2$.
\end{stwierdzenie}

\begin{proof}
Wstawiamy nowy wierzchołek $x$ do wnętrza trójkąta $abc$.
Po ruchu $1\to3$ trójkąty tej części to $abx$, $bcx$, $cax$;
trójkąt $abd$ pozostaje bez zmian. Wspólna krawędź $ab$
trójkątów $abx$ i $abd$ ma teraz przeciwległe wierzchołki
$x,d$. Krawędź $xd$ jeszcze nie istnieje, ponieważ $x$ jest
nowy, a jego sąsiadami są dotąd $a,b,c$; cztery wierzchołki
$ a,b,x,d$ są różne. Ruch $2\to2$ jest więc dopuszczalny.
Zamieniamy przekątną $ab$ na $xd$. Powstają trójkąty $adx$, $bdx$, $bcx$, $cax$.
Są one dokładnie czterema trójkątami otrzymanymi przez
podział dwóch starych trójkątów wzdłuż krawędzi $ab$.
Zależność od położenia geometrycznego $x$ jest pozorna:
liczy się kombinatoryka kompleksu, a oba podziały
dysku są homeomorficzne względnie do jego brzegu.
\end{proof}

\begin{twierdzenie}[Pachner; wynik przyjmowany]
Niech $K,L$ będą dwiema skończonymi triangulacjami
tej samej zwartej, zamkniętej rozmaitości PL
wymiaru dwa, z ustalonymi triangulacjami zgodnymi
z jej strukturą PL.
Wtedy istnieje skończony ciąg dopuszczalnych ruchów
$1\leftrightarrow3$ i $2\leftrightarrow2$ prowadzący
od $K$ do kompleksu izomorficznego z $L$.
\end{twierdzenie}

\begin{uwaga}[Status dowodu]
Udowodniliśmy lokalną poprawność ruchów i ich wpływ na
$\chi$, ale twierdzenie o istnieniu ciągu łączącego
\emph{dowolne} dwie triangulacje jest znacznie silniejsze.
Wymaga teorii podziałów gwiaździstych i wspólnych
podziałów PL (lub pełnego argumentu Pachnera).
Przyjmujemy je jako klasyczny wynik Pachnera
(U. Pachner, \emph{P.L. homeomorphic manifolds are
equivalent by elementary shellings}, European Journal of
Combinatorics, 1991); pełnego dowodu tutaj nie podajemy.
\end{uwaga}

